\chapter{Design SOLID}
\section{Archittetura e Design}
Ogni software rappresenta due valori principali per l'investitore:
\begin{itemize}
    \item \textbf{\textcolor{brown}{Comportamento}}, cioè le funzioni principali che svolge.
    \item \textbf{\textcolor{brown}{Struttura}}, cioè un programma facilmente modificabile, perché verrà modificato se è utilizzato.
\end{itemize}

Una volta che i requisiti sono chiari, si trasformano in \textbf{design} che ne definisce i componenti e le loro interazioni ed è diviso in due fasi:
\begin{enumerate}
    \item \textbf{FASE ARCHITETTURALE:} visione di alto livello, proprietà esterne e relazioni, dettate dai requisiti funzionali e non.
    \item \textbf{FASE DETTAGLIO DESIGN:} decomposizione dei componenti in dettagli, dall'architettura ed i requisiti funzionali.
\end{enumerate}

\section{Principi SOLID}
\begin{itemize}
    \item \textbf{\textcolor{brown}{Single Responsibility Principle (SRP)}}, la classe deve avere un solo motivo per cambiare:
    \begin{itemize}
        \item Deve avere una sola funzionalità, per evitare di creare un malfunzionamento in caso di cambiamento.
        \item È difficile da implementare perché le classi diventano tante e perché siamo abituati a raggruppare le funzionalità.
        \item Un modulo è un file sorgente che coinvolge più classi ad alta coesione che rispettano il principio \textbf{SRP}.
    \end{itemize}
    \textbf{\textcolor{red}{SOLUZIONE:}} ogni classe deve contenere il codice per la sola funzione che gli spetta ed utilizza il pattern FACADE.
    \item \textbf{\textcolor{brown}{Open Close Principle (OCP)}}, deve essere possibile estendere il comportamento di una classe senza modificarla:
    \begin{itemize}
        \item Ogni modifica su codice già rilasciato è vietata, bisogna creare un nuovo metodo o meglio una nuova classe.
        \item Migliora la stabilità, manutenzione e riutilizzabilità del codice; aperta per implementazioni, chiusa per modifiche
        \item Le variabili sono private, non si usano globali. Le frecce piene significano implementazioni, vuote invece relazioni. Se una freccia va dalla classe A alla B, significa che all'interno della classe A viene menzionata e usata la classe B.
        \item In pratica utilizzo un'interfaccia ed una o più implementazioni con polimorfismo per definire un modulo ad alto livello, protetto.
    \end{itemize}
    \textbf{\textcolor{red}{SOLUZIONE:}} l'astrazione creando un interfaccia, i cui oggetti istanziati saranno attributi della classe originaria scelti con i setter.
    \item \textbf{\textcolor{brown}{Liskov Substitution Principle (LSP):}}, la classi derivate devono essere sostitutive della loro classe base senza problemi:
    \begin{itemize}
        \item La validità di un modello dipende dall'utilizzatore, dato che il focus è il comportamento pubblico della classe OOP.
    \end{itemize}
    \textbf{\textcolor{red}{SOLUZIONE:}} delegare i metodi alle sottoclassi che ereditano da quella base che contiene le firme delle funzioni.
    \item \textbf{\textcolor{brown}{Interface Segregation Principle (ISP)}}, creare gerarchie di interfacce, in modo che il client non sia costretto a dipendere da metodi non usati:
    \begin{itemize}
        \item Evitare che si utilizzino interfacce troppo ricche di metodi, perché potrebbero essere inutili per molte sottoclassi.
        \item In caso di modifiche ricompilo solo le interfacce modificate, analogamente riduco le dipendenze alle essenziali
    \end{itemize}
    \textbf{\textcolor{red}{SOLUZIONE:}} creo più interfacce legate tra di loro, ognuna che specifica maggiormente la precedente, non una.
    \item \textbf{\textcolor{brown}{Dependency Inversion Principle (DIP)}}, dipendere da astrazioni (stabili e non modificabili), non da classi specifiche:
    \begin{itemize}
        \item Se un modulo di alto livello dipende da uno di basso livello non va bene! (esempio di pizza/pizzeria).
        \item Deve dipendere da un interfaccia, non da una concretizzazione della stessa.
        \item \textbf{\textcolor{blue}{Dependency Injection}}, detta anche inversione di controllo, configura le dipendenze dall'esterno non staticamente. Iniettare tramite parametro nel costruttore, tramite un metodo setter esposto o tramite interfaccia con un iniettore.
    \end{itemize}
     \textbf{\textcolor{red}{SOLUZIONE:}} introdurre un livello di astrazione tra alto e basso livello, che dipendono da un'interfaccia comune.
\end{itemize}

\section{Principi di coesione dei componenti}
I primi due sono inclusivi, mentre il terzo esclusivo e sono tutti essenziali per evitare sbilanciamenti. Un componente è l'unità di codice più piccola che può essere distribuita e compilata indipendentemente, che devono rimanere coesi quando scelgo le classi:
\begin{itemize}
    \item \textbf{\textcolor{brown}{Reuse/Release Equivalence (REP):}}rilascio associato ad un numero e comporta la coerenza dei componenti nei pacchetti che vengono distribuiti.
    \item \textbf{\textcolor{brown}{Common Closure Principle (CCP):}} raccolgo in componenti le classi che possono cambiare per la stessa ragione, analogia con il principio solid \textbf{SRP}.
    \item \textbf{\textcolor{brown}{Common Reuse Principle (CRP):}} non costringere l'utilizzatore a dipendere da codice di cui non ha bisogno, analogia con \textbf{ISP}.
\end{itemize}

\section{Principi di accoppiamento dei componenti}
Ci permettono di gestire le relazioni tra i componenti:
\begin{itemize}
    \item \textbf{\textcolor{brown}{Acyclic Dependencies Principle (ADP)}}, impedire cicli nel grafo delle dipendenze perché se si lavora in team si rischia di rompere il lavoro degli altri.
    \item \textbf{\textcolor{brown}{Stable Dependencies Principle (SDP)}}, si cerca di dipendere nella direzione della stabilità ed è importante capire i componenti stabili e instabili.
    \item \textbf{\textcolor{brown}{Stable Abstraction Principle (SAP)}}, un componente deve far dipendere il suo livello di astrazione dal suo livello di stabilità.
\end{itemize}